% !TEX TS-program = pdflatex
% THESIS_OUTLINE.tex
% Chapter-by-chapter outline of the REE 599 final thesis, with a map from
% every existing project artifact to its destination chapter.
% Companion document: FIVE_DAY_WRAPUP_PLAN.tex (schedule for closing the
% remaining work). Format follows PROJECT_COMPLETION_PLAN.tex.
\documentclass[11pt,a4paper]{article}

\usepackage[margin=2.3cm]{geometry}
\usepackage[T1]{fontenc}
\usepackage{lmodern}
\usepackage{microtype}
\usepackage{setspace}
\onehalfspacing
\usepackage{amsmath,amssymb}
\usepackage{booktabs}
\usepackage{longtable}
\usepackage{colortbl}
\usepackage[table]{xcolor}
% Compat shim: colortbl 1.0k (2024) expects \insert@pcolumn from array v2.6,
% but this MiKTeX tree carries array v2.5g. Without the alias every p-column
% in a colortbl document throws "Undefined control sequence".
\makeatletter
\@ifundefined{insert@pcolumn}{\let\insert@pcolumn\insert@column}{}
\makeatother
\usepackage{enumitem}
\usepackage[hidelinks]{hyperref}
\usepackage{titlesec}
\titleformat{\section}{\large\bfseries}{\thesection.}{0.5em}{}
\titleformat{\subsection}{\normalsize\bfseries}{\thesubsection}{0.5em}{}
\emergencystretch=1.5em   % long typewriter file names need slack

\definecolor{rowExist}{RGB}{198,239,206}   % material exists, edit/merge only
\definecolor{rowPart}{RGB}{255,235,156}    % partly exists, needs new writing
\definecolor{rowNew}{RGB}{252,213,180}     % must be written new / test-dependent

% File-name macro: typewriter face with line breaks allowed after underscores,
% so long artifact names do not overflow the margin.
\newcommand{\f}[1]{{\ttfamily\def\_{\textunderscore\allowbreak}#1}}

\title{\textbf{REE 599 Final Thesis: Chapter Outline and Source Map}\\[0.4em]
\large Autonomous Thermoelectric Atmospheric-Water PV-Cleaning Robot\\[0.2em]
\normalsize for Egypt's New Sustainable Cities (Commercial / Residential / Street-Lamp PV)}
\author{Rahma El-Hamouly, Graduation Project}
\date{Issued: 10 June 2026}

\begin{document}
\maketitle

\begin{abstract}
\noindent
This document is the writing plan for the REE 599 thesis (target 80--100 pages).
For each chapter it states the purpose, the section list, which existing project
files supply the content, and what still has to be written from scratch or
filled in after the 4-tile prototype test. The structure follows the REE 598
departmental guidelines (\f{REE 598 GP\_1\_Guidelines.pdf}): cover page,
acknowledgment, one-page abstract, table of contents, lists of figures and
tables, nomenclature, numbered chapters, references, and appendices, with
chapter-based numbering of sections, figures, tables, and equations
(Times~12\,pt, 1.5 line spacing). The day-by-day schedule that executes this
outline is in \f{FIVE\_DAY\_WRAPUP\_PLAN.tex}.
\end{abstract}

\tableofcontents
\newpage

% =============================================================================
\section{Ground Rules}
% =============================================================================
\begin{enumerate}[leftmargin=*]
  \item \textbf{Single thesis source file.} The thesis is assembled as
        \f{REE599\_thesis.tex} (\f{report} class), reusing the siunitx shim and
        preamble of \f{prototype\_design\_report.tex}, which compiles cleanly on
        this TeX tree.
  \item \textbf{Numbers come from the verified engine.} Every performance figure
        traces to \f{scripts/tec\_robot\_design.py} and its archived outputs
        (\f{output/tec\_robot\_results.json}, \f{output/tec\_robot\_monthly.csv}).
        If a parameter changes during design closure (heat-sink resistance,
        battery size), the script is re-run before the text is updated.
  \item \textbf{The eight review flags are resolved before assembly.} The open
        items in \f{docs/REE598\_PROPOSAL\_REVIEW.md} \S3.1--3.8 contradict
        numbers that would otherwise be copied into Chapters 5, 7, and 11.
        Resolving them first (Day 1 of the wrap-up plan) avoids writing the same
        chapter twice.
  \item \textbf{Citations.} IEEE style via \f{IEEEtran}/\f{ieeetr} and
        \f{references.bib} (about 34 entries, all with DOI/URL). The citation
        corrections logged in \f{docs/CITATION\_VERIFICATION\_2026-06-02.md}
        carry over: the 0.3~L/m$^2$ microfibre cleaning-water dose stays a
        declared design assumption (within the published 0.1--0.5~L/m$^2$ band
        for water-assisted robotic PV cleaning), not a literature value.
  \item \textbf{Benban is evidence, not the deliverable.} Per
        \f{memory/project\_context.md}, the Benban work appears once, in
        Chapter~3, as climate-feasibility validation. The robot design carries
        the rest of the thesis.
\end{enumerate}

% =============================================================================
\section{Front Matter (Guideline Items 1--6)}
% =============================================================================
\begin{itemize}[nosep,leftmargin=*]
  \item Cover page: faculty template with \f{logo.jpg}, official title,
        supervisor, term. Reuse the title page of \f{prototype\_design\_report.tex}
        as the starting layout.
  \item Acknowledgment: one page, written new (supervisor, department, vendors
        who supplied datasheets).
  \item Abstract: one page. Written new, last, after measured prototype results
        exist. Must quote the final headline numbers: per-tile yield, 69-tile
        sizing, 1.011~L/day annual mean, COP 1.44, measured prototype yield, and
        the feasibility headline (at the residential 0.3~L/m$^2$ dose, full water
        self-supply is feasible from a grid-coupled TEC array, about 14\% of the
        host array's annual energy).
  \item Table of contents, list of figures, list of tables: generated.
  \item Nomenclature and abbreviations: written new, one page. Harvest symbols
        from \f{prototype\_design\_report.tex} (TEC coefficients $\alpha, R, K$;
        psychrometric symbols $\omega, T_{dew}, P_{sat}$; energy symbols) and
        abbreviations (AWG, AWH, TEC, VCR, SEC, COP, LCOW, PSH, PR, DoD, BOM).
\end{itemize}

% =============================================================================
\section{Chapter-by-Chapter Outline}
% =============================================================================

% -----------------------------------------------------------------------------
\subsection{Chapter 1: Introduction (5 pp)}
\textit{Purpose.} State the soiling/water-scarcity problem in arid Egypt, present
the robot-with-self-generated-water concept, and define scope across the three
deployment scales tied to the new sustainable cities.

\textbf{Sections:} 1.1 Importance of the field (PV growth in Egypt, soiling
losses 20--30\%/yr untreated); 1.2 The water-logistics problem of wet cleaning;
1.3 Proposed solution: autonomous TEC-AWG cleaning robot; 1.4 Deployment scales
A/B/C and target cities; 1.5 Contributions; 1.6 Thesis structure.

\textbf{Existing material:} \f{main\_report.tex} ch.~1 (motivation, proposed
solution, scope, contribution) merges with \f{REE598\_proposal.tex} ch.~1 and
\f{prototype\_design\_report.tex} ch.~1 (problem statement, functional
requirements table). \f{memory/project\_context.md} fixes the scale ordering:
street-lamp and rooftop scales lead, industrial is the stretch case (review
flag \S4.1).

\textbf{Write new:} contributions list (now includes measured prototype data);
thesis-structure roadmap paragraph.

% -----------------------------------------------------------------------------
\subsection{Chapter 2: Literature Review (10 pp)}
\textit{Purpose.} Survey soiling in MENA, wet vs.\ dry cleaning, AWH
technologies, TEC condensers in arid climates, and existing cleaning robots,
ending with the research gap the project fills.

\textbf{Sections:} 2.1 PV soiling in Egypt and MENA; 2.2 Cleaning methods and
why wet cleaning wins on recovery; 2.3 Atmospheric water harvesting technologies
(sorbent, fog, VCR, TEC); 2.4 TEC condenser physics and prior arid-climate
implementations; 2.5 Autonomous PV-cleaning robots, state of the art; 2.6
Egyptian new cities and solar context; 2.7 Critique and research gap.

\textbf{Existing material:} \f{REE598\_proposal.tex} ch.~2 (five sections,
already cited) and \f{main\_report.tex} ch.~2 cover all of 2.1--2.6. The AWH
comparison framework in \f{project\_analysis\_awh\_comparison\_and\_scope.tex}
supplies the SEC bands for 2.3.

\textbf{Write new:} merge pass only; deduplicate the two existing reviews and
sharpen 2.7 so the gap statement matches the pivoted deliverable (no published
robot generates its own cleaning water on-board).

% -----------------------------------------------------------------------------
\subsection{Chapter 3: Feasibility Evidence: TEC Water Harvesting in the
Benban Climate (8 pp)}
\textit{Purpose.} Compress the Benban-era analyses into the evidence base: TEC
beats VCR on sustainability and maintenance at small scale, and the Benban
hourly climate proves condensation is achievable in arid Egypt.

\textbf{Sections:} 3.1 Benban as reference climate dataset (NASA POWER, 8784
hourly records, 2020); 3.2 Psychrometric analysis and the seasonal condensation
window; 3.3 TEC vs.\ VCR: efficiency, sustainability, maintenance; 3.4 Scaling
floor: where TEC wins (small/distributed) and where VCR wins (central plant);
3.5 Egyptian PV inventory and the addressable small-scale segment; 3.6
Conclusion: design basis for a TEC robot.

\textbf{Existing material:} \f{BENBAN\_TEC\_VCR\_CAPACITY\_ANALYSIS.tex};
\f{tec\_vcr\_sustainability\_efficiency\_maintenance\_v1.2.tex} (efficiency,
sustainability, maintenance, lifecycle sections);
\f{tec\_single\_vs\_parallel\_deployment\_analysis\_v1.0.tex} and
\f{tec\_deployment\_analysis\_v1.1\_egypt\_market\_supplement.tex} (deployment
configurations, Egyptian market);
\f{egypt\_pv\_inventory\_and\_tec\_small\_scale.tex} (inventory, scaling floor,
recommendation matrix); \f{main\_report.tex} ch.~3 (climate analysis,
validation against measurements); \f{output/psychrometrics\_monthly.csv}.

\textbf{Write new:} compression only. These five standalone documents total
well over 60 pages; the chapter keeps the summary tables and one psychrometric
figure and cites the standalone PDFs as project archive material in an
appendix pointer.

% -----------------------------------------------------------------------------
\subsection{Chapter 4: System Architecture and Deployment Scales (8 pp)}
\textit{Purpose.} Define the robot as a system (AWG + cleaning + mobility +
power + control), specify the three scale variants, and record the architecture
decisions forced by the feasibility finding.

\textbf{Sections:} 4.1 Overall concept and subsystem block diagram; 4.2 Scale A,
street-lamp pole unit; 4.3 Scale B, commercial rooftop rail-guided carriage
(prototype scale); 4.4 Scale C, industrial: fleet of Scale-B units, not a gantry;
4.5 Per-array TEC harvester sizing and reservoir buffering; 4.6 Docking and
safe-park station; 4.7 Requirements tables per scale.

\textbf{Existing material:} \f{main\_report.tex} ch.~4 (overall concept, design
requirements per scale); \f{REE598\_proposal.tex} ch.~3 (solution architecture,
deployment-scale specification tables); \f{figures/system\_architecture.svg};
\f{prototype\_design\_report.tex} feasibility section (per-array harvester,
full water self-supply at the 0.3~L/m$^2$ dose).

\textbf{Write new:} 4.4 must be rewritten to the fleet decomposition that
resolves review flag \S3.7 (the 1.67-tonne gantry reservoir is dropped); 4.5
must state the multi-night harvest-to-buffer sequence that resolves flag \S3.8
(condensate tank sizing); 4.6 docking detail is currently a one-line mention
and needs a figure plus a half page of mechanism description (charging contacts,
park orientation, wind/dust survival).

% -----------------------------------------------------------------------------
\subsection{Chapter 5: TEC AWG Design: Psychrometrics and Operating Point (12 pp)}
\textit{Purpose.} The thermodynamic core: module selection, device coefficients,
the coupled heat-and-mass-transfer condensation model, the operating-point
solver, and tile-array sizing.

\textbf{Sections:} 5.1 TEC1-12706 selection and Egyptian-market coefficients;
5.2 Governing equations ($Q_C$, $W$, COP); 5.3 Site psychrometrics (Buck,
Magnus, humidity ratio); 5.4 Condensation as coupled heat/mass transfer (Lewis
analogy); 5.5 Hot-side thermal path and the self-consistent operating point;
5.6 Operating modes A/B and the smart night-only strategy; 5.7 Condensation
margin sensitivity; 5.8 Tile count: 69 tiles for 1~L/day; 5.9 SEC and
reconciliation with published TEC-AWG data.

\textbf{Existing material:} \f{prototype\_design\_report.tex} chs.~2--4 carry
almost the whole chapter, including the strategy and margin-sensitivity tables
and \f{figures/fig\_margin\_sensitivity.pdf}; \f{main\_report.tex} ch.~5
duplicates parts and is merged in; \f{output/tec\_robot\_results.json} for all
numbers; \f{scripts/verify\_independent\_solver.py} for the solver
cross-check paragraph.

\textbf{Write new:} 5.5 and 5.9 absorb review flags \S3.2--3.4. The proposal's
SEC of 0.927~kWh/L is superseded by the coupled-solver value (4.16~kWh/L
condensation-only, 5.27 with parasitics), which sits inside the published
2.8--12~kWh/L band; the chapter must say so explicitly and show the heat-sink
resistance sensitivity (design value vs.\ realistic 0.5--1.0~K/W per module
pair). Section 5.9 gains a measured point from the prototype
(test-dependent).

% -----------------------------------------------------------------------------
\subsection{Chapter 6: Robot Mechanical and Cleaning Design (8 pp)}
\textit{Purpose.} Mechanical design of the rail-guided carriage: drive sizing on a
tilted panel, the fall-line sweep, cleaning head, water circuit, and the CAD
assembly.

\textbf{Sections:} 6.1 Chassis and mass budget (8~kg rail-guided carriage); 6.2
Drive-system sizing (42.7~N resistive force, rack-and-pinion at $\geq$0.84~N$\cdot$m
on a 30~mm pinion near 38~rpm, 1.5 safety factor) and why the free
differential-drive crawler was revised to a rail carriage (it cannot hold yaw on
24$^\circ$ glass, confirmed in Gazebo); 6.3 Fall-line sweep coverage (full-width
brush, one down-and-up traverse, 66.7~s per 2~m$^2$ panel); 6.4 Cleaning
head: full-width 0.96~m brush plus low-pressure spray, 0.3~L/m$^2$ microfibre dose,
102.5 events/yr; 6.5 Water circuit: gutter, reservoir, pump, level sensing; 6.6 CAD
assembly model and drawings; 6.7 Structural and safety considerations (end-of-travel
detection, tilt trip).

\textbf{Existing material:} \f{prototype\_design\_report.tex} ch.~6 (drive
sizing, coverage, water logic); \f{main\_report.tex} ch.~6 (per-scale
mechanical sections, structural/safety) and ch.~7 (water storage, treatment,
delivery) folds into 6.5.

\textbf{Write new:} 6.6 does not exist. A block-level CAD assembly (Fusion 360)
with one exploded view and overall-dimension drawing is produced on Day~4 of
the wrap-up plan; renders go here, working drawings go to Appendix~D.

% -----------------------------------------------------------------------------
\subsection{Chapter 7: Power and Energy System (6 pp)}
\textit{Purpose.} Energy budget, PV and battery sizing, and the corrected
Scale-B electrical BOM.

\textbf{Sections:} 7.1 Daily energy budget (5.32~kWh/day, TEC stack 79\%);
7.2 PV sizing (1147~Wp, 5.34~m$^2$, PSH 5.8, PR 0.80); 7.3 Battery sizing
(698~Ah at 12~V, 80\% DoD) and the corrected Scale-B pack (2.5~kWh, realistic
450~USD pricing); 7.4 DC architecture, MPPT, converters, protection; 7.5 Water
self-supply feasibility: at the 0.3~L/m$^2$ dose a grid-coupled TEC array meets the
full residential demand (about 38 tiles, roughly 14\% of the host array's annual
energy), so the design consequence is sizing the harvester, not rationing water.

\textbf{Existing material:} \f{prototype\_design\_report.tex} ch.~5 (budget,
sizing, feasibility finding) and \f{figures/fig\_energy\_balance.pdf};
\f{main\_report.tex} ch.~8 (charge controller, DC architecture);
\f{REE598\_proposal.tex} \S6.4.

\textbf{Write new:} 7.3 incorporates review flags \S3.5--3.6 (battery
undersizing and unrealistic pack price); the Scale-B cost figure changes to
about 900~USD and propagates to Chapter~11 economics. 7.4 needs one wiring
one-liner per protection element (fuse, reverse-polarity, low-voltage cutoff),
currently scattered.

% -----------------------------------------------------------------------------
\subsection{Chapter 8: Control, Autonomy, and Firmware (6 pp)}
\textit{Purpose.} The supervisory state machine, the cold-plate PID loop, safety
logic, and the firmware as actually implemented on the prototype ESP32.

\textbf{Sections:} 8.1 Supervisory state machine
(\textsc{Idle}/\textsc{Harvest}/\textsc{Store}/\textsc{Clean}/\textsc{Dock});
8.2 Dew-point computation on-board; 8.3 Cold-plate set-point control, PID with
anti-windup at $I_{max}$; 8.4 Safety logic (over-current, limit switches, IMU
trip, watchdog); 8.5 Implemented firmware: architecture, set-point hold, and
logging (ESP32, DHT22, DS18B20, INA219); 8.6 Observed control performance.

\textbf{Existing material:} \f{prototype\_design\_report.tex} ch.~7 covers
8.1--8.4 as design (pseudocode level); \f{main\_report.tex} ch.~8 control
section overlaps.

\textbf{Write new:} 8.5 and 8.6 depend on the prototype build (Days 2--3) and
the test run (Day 4). The firmware listing itself goes to Appendix~C; the
chapter shows the loop timing, the measured set-point step response, and the
hold accuracy.

% -----------------------------------------------------------------------------
\subsection{Chapter 9: Annual Simulation Results (8 pp)}
\textit{Purpose.} Present the hourly Benban-year simulation of the full 69-tile
system: monthly water yield, COP, energy, and the seasonal pattern that drives
the deployment concept.

\textbf{Sections:} 9.1 Method (8784 hourly records, night-hours window,
per-tile linear scaling); 9.2 Climate inputs; 9.3 Monthly yield: 4.5~L/day
November peak, near zero June--September, 1.011~L/day annual mean; 9.4 COP and
SEC vs.\ temperature lift; 9.5 Smart vs.\ all-night strategy comparison; 9.6
Implications: reservoir banking of the October--March surplus.

\textbf{Existing material:} \f{prototype\_design\_report.tex} ch.~8 (method,
monthly table, discussion); \f{output/tec\_robot\_monthly.csv} and
\f{output/tec\_robot\_results.json} (all numbers);
\f{figures/fig\_monthly\_climate.pdf}, \f{fig\_diurnal.pdf},
\f{fig\_monthly\_water.pdf}, \f{fig\_cop\_sec\_dT.pdf};
\f{main\_report.tex} ch.~10 results section.

\textbf{Write new:} nothing structural. If design closure changes the heat-sink
resistance (flag \S3.4), the simulation is re-run and the tables regenerate
from the CSV; the text already reads off the script outputs.

% -----------------------------------------------------------------------------
\subsection{Chapter 10: Prototype Build and Experimental Validation (10 pp)}
\textit{Purpose.} The 4-tile bench demonstrator: rationale, BOM, build, test
protocol, measured results, uncertainty, and model comparison. This is the
chapter that distinguishes REE 599 from REE 598.

\textbf{Sections:} 10.1 Rationale: per-tile linearity, 4 tiles predict
59~g/day at the Benban annual mean; 10.2 Bill of materials, Egyptian market
($\approx$5{,}760~EGP full kit; bench-critical subset $\approx$3{,}800~EGP);
10.3 Assembly (TEC sandwich, cold plate, gutter, shared cold-side fan); 10.4
Instrumentation and calibration (DHT22, DS18B20 ice-point check, INA219); 10.5
Test protocol (overnight runs, hourly or end-of-run condensate weighing); 10.6
June test conditions: humidified-enclosure method and dawn-hour outdoor runs,
and why point-wise model validation at measured $(T, RH)$ replaces a direct
comparison with the annual-mean 59~g/day figure; 10.7 Measured results: yield,
COP, set-point hold; 10.8 Uncertainty analysis; 10.9 Model vs.\ measurement and
what it means for the 69-tile scale-up.

\textbf{Existing material:} \f{prototype\_design\_report.tex} prototype and
test-plan chapters carry 10.1--10.5 and the uncertainty formulation almost
verbatim.

\textbf{Write new:} 10.6--10.9 are entirely test-dependent (wrap-up plan Days
3--5). Data tables, the predicted-vs-measured plot, and photographs of the rig
are new artifacts. This is the largest block of genuinely new thesis text.

% -----------------------------------------------------------------------------
\subsection{Chapter 11: Deployment in Egyptian New Cities and Economics (8 pp)}
\textit{Purpose.} Translate the design into a rollout plan for the New
Administrative Capital, New Alamein, New Mansoura, and New Ismailia, with
corrected per-scale economics.

\textbf{Sections:} 11.1 Target building stock and PV inventory per city; 11.2
Scale-matched deployment (A: street-light grids; B: commercial rooftops; C:
fleets on industrial yards); 11.3 Corrected economics: energy recovered from
the consistent soiling model, realistic battery pricing, payback per scale;
11.4 Water-economics comparison vs.\ trucked water; 11.5 Environmental, social,
and regulatory considerations; 11.6 Rollout phasing.

\textbf{Existing material:} \f{egypt\_pv\_inventory\_and\_tec\_small\_scale.tex}
(inventory, recommendation matrix);
\f{tec\_deployment\_analysis\_v1.1\_egypt\_market\_supplement.tex} (market
pricing); \f{REE598\_proposal.tex} ch.~6 economics tables and ch.~5
(regulations, standards, ethical/environmental sections);
\f{main\_report.tex} ch.~10 economic analysis.

\textbf{Write new:} 11.3 is a rebuild, not a copy. Review flag \S3.1 showed the
proposal's Table~6.2 overstates recovered energy by 3.5--4.7$\times$; the
corrected Scale-B case is roughly 3{,}360~kWh/yr recovered, 67~USD/yr at the
0.02~USD/kWh tariff, and payback near 13 years at the corrected 900~USD unit
cost. The honest framing (payback driven by labour and water-logistics savings,
not tariff revenue) must be written fresh. 11.6 is new but short.

% -----------------------------------------------------------------------------
\subsection{Chapter 12: Conclusions and Future Work (3 pp)}
\textit{Purpose.} State the validated findings and the recommendations.

\textbf{Sections:} 12.1 Conclusions (per-tile physics validated by experiment;
69-tile sizing; seasonal yield pattern; water self-supply feasible at the
0.3~L/m$^2$ dose; corrected economics); 12.2 Limitations (single climate year,
bench-scale
validation, June-only test window); 12.3 Future work (field trial on a real
rooftop array, winter validation campaign, docking-station prototype, fleet
coordination, MOF-assisted hybrid harvesting).

\textbf{Existing material:} \f{prototype\_design\_report.tex} conclusions
chapter; \f{main\_report.tex} ch.~11.

\textbf{Write new:} rewritten last, after measured data lands; half a day.

% =============================================================================
\section{Appendices}
% =============================================================================
\begin{itemize}[nosep,leftmargin=*]
  \item \textbf{Appendix A: Simulation code.} \f{scripts/tec\_robot\_design.py}
        and \f{scripts/verify\_independent\_solver.py} listings with a usage
        note; archived outputs \f{tec\_robot\_results.json},
        \f{tec\_robot\_monthly.csv}, \f{tec\_margin\_sensitivity.csv}.
  \item \textbf{Appendix B: Bills of materials.} B.1 the 4-tile bench BOM
        (exists, \f{prototype\_design\_report.tex}); B.2 the full-scale 69-tile
        Scale-B BOM with Egyptian vendor pricing (new, produced Day~4 of the
        wrap-up plan); B.3 corrected Scale-B electrical BOM
        (battery upsized per flags \S3.5--3.6).
  \item \textbf{Appendix C: Firmware listing.} ESP32 source as built
        (new, Days 2--3).
  \item \textbf{Appendix D: CAD drawings.} Assembly drawing, exploded view,
        overall dimensions (new, Day~4).
  \item \textbf{Appendix E: Test data.} Raw logs (CSV), condensate weighing
        sheets, calibration records (new, Days 3--5).
  \item \textbf{Appendix F: Benban climate data.} Monthly summary from
        \f{output/psychrometrics\_monthly.csv} and pointers to the archived
        standalone Benban analyses (existing PDFs).
  \item \textbf{Appendix G: Standards and regulations.} Carried from
        \f{REE598\_proposal.tex} Appendix D.
\end{itemize}

% =============================================================================
\section{Source-to-Chapter Map (Summary)}
% =============================================================================
Row colour: green = exists, merge/edit only; yellow = exists in part, needs new
sections; orange = mostly new or test-dependent.

\begin{center}
\small
\renewcommand{\arraystretch}{1.25}
\begin{longtable}{@{}clp{7.2cm}c@{}}
\toprule
\textbf{Ch.} & \textbf{Title (short)} & \textbf{Primary sources} & \textbf{Pages} \\
\midrule
\endhead
\rowcolor{rowExist} 1 & Introduction & \f{main\_report} ch.1; \f{REE598\_proposal} ch.1; \f{prototype\_design\_report} ch.1 & 5 \\
\rowcolor{rowExist} 2 & Literature review & \f{REE598\_proposal} ch.2; \f{main\_report} ch.2; \f{project\_analysis\_awh\_comparison} & 10 \\
\rowcolor{rowExist} 3 & Feasibility evidence (Benban) & \f{BENBAN\_TEC\_VCR\_CAPACITY\_ANALYSIS}; \f{tec\_vcr\_...\_v1.2}; \f{tec\_...\_v1.0/v1.1}; \f{egypt\_pv\_inventory}; \f{main\_report} ch.3 & 8 \\
\rowcolor{rowPart} 4 & Architecture, three scales & \f{main\_report} ch.4; \f{REE598\_proposal} ch.3; flags \S3.7--3.8 resolutions & 8 \\
\rowcolor{rowPart} 5 & TEC AWG design & \f{prototype\_design\_report} chs.2--4; \f{tec\_robot\_results.json}; flags \S3.2--3.4 & 12 \\
\rowcolor{rowPart} 6 & Mechanics and cleaning & \f{prototype\_design\_report} ch.6; \f{main\_report} chs.6--7; CAD (new) & 8 \\
\rowcolor{rowPart} 7 & Power and energy & \f{prototype\_design\_report} ch.5; \f{main\_report} ch.8; flags \S3.5--3.6 & 6 \\
\rowcolor{rowPart} 8 & Control and firmware & \f{prototype\_design\_report} ch.7; firmware as built (new) & 6 \\
\rowcolor{rowExist} 9 & Annual simulation & \f{prototype\_design\_report} ch.8; \f{tec\_robot\_monthly.csv}; 6 figures & 8 \\
\rowcolor{rowNew} 10 & Prototype and validation & \f{prototype\_design\_report} build/test chapters + measured data (new) & 10 \\
\rowcolor{rowNew} 11 & Deployment and economics & \f{egypt\_pv\_inventory}; \f{tec\_deployment\_v1.1}; corrected economics (new) & 8 \\
\rowcolor{rowNew} 12 & Conclusions, future work & \f{prototype\_design\_report} conclusions + measured findings & 3 \\
 & Front matter, references & guidelines items 1--6; \f{references.bib} & 10 \\
 & Appendices A--G & scripts, BOMs, firmware, CAD, test data, climate, standards & 8--12 \\
\midrule
 & \textbf{Total} & & \textbf{$\approx$100} \\
\bottomrule
\end{longtable}
\end{center}

If the total runs over 100 pages, Chapter 3 compresses first (its standalone
PDFs remain in the archive) and Appendix~A switches from full listings to
repository pointers.

% =============================================================================
\section{Items Blocked on Prototype Results}
% =============================================================================
Everything else can be assembled from existing files in parallel with the
build. Only these wait for measured data:
\begin{enumerate}[nosep,leftmargin=*]
  \item Chapter 10, \S\S10.6--10.9 (measured yield, COP, uncertainty,
        model comparison).
  \item Chapter 8, \S\S8.5--8.6 (firmware as built, set-point hold performance).
  \item Chapter 5, \S5.9 measured SEC point.
  \item Chapter 12 conclusions and the abstract headline numbers.
  \item Appendices C and E.
\end{enumerate}

\end{document}
